\paragraph{Statement A.5.8 (Nonvanishing and both square-root oscillation signs).}
For the actual sharp arc $I$, there are infinitely many positive
spectral parameters $t$ such that
\[
 \mathcal C(t)=\sum_{j:t_j=t}p'_jQ_{j,I}\ne0.
\]
Consequently
\[
 R(X)-C_0X=\Omega_+(X^{1/2})
 \quad\hbox{and}\quad
 R(X)-C_0X=\Omega_-(X^{1/2}).
\]
In particular the requested unsigned $\Omega(X^{1/2})$ assertion holds.

\begin{proof}
We prove nonvanishing by a localized positive-frequency relative-wave
calculation, including its endpoint contribution. The oscillation
conclusion then follows from Statement A.5.6.
Throughout this proof set
\[
 a=17,\quad B=12\sqrt2,\quad R=\operatorname{arcosh}(17),
 \quad d_0=\operatorname{arcosh}(3),\quad
 \tau(u)=\operatorname{arsinh}(B\cosh u).
\]
Thus $R=2d_0$.
The period normalization throughout is
$p'_j=\sum_{i=1}^2\int_0^{L_i}u_j(\ell_i(v))\,dv$,
with the primitive cyclic translation stabilizers $H_i$ used in
A.3.1. If a larger axis normalizer identifies the two orientations,
unfolding with $H_i$ retains the resulting positive duplicate
multiplicity; it does not divide a period by an unrecorded factor.

\emph{1. The frequency window and the exact source functional.}
Choose a real even Schwartz function $\rho$ whose Fourier transform
\[
 \widehat\rho(y)=\int_{\mathbb R}\rho(t)e^{-ity}\,dt
\]
is smooth, nonnegative, supported in $(-\varepsilon,\varepsilon)$,
and positive at zero. The sufficiently small fixed $\varepsilon>0$
will be specified below.
Consider the absolutely convergent sum
\begin{equation}\label{A:nonvanishing-window}
 \mathscr S_T=
 \sum_{j\geq1}p'_jQ_{j,I}\rho(t_j-T)e^{it_jd_0},
 \qquad T\to+\infty.
\end{equation}
The polynomial individual bounds already proved in Statement A.5.6
justify this convergence. We shall prove
\begin{equation}\label{A:nonvanishing-wave-asymptotic}
 \mathscr S_T=C_*e^{-i\pi/4}T^{-1/2}+O(T^{-1}),
 \qquad C_*>0.
\end{equation}
If every complete-eigenspace contraction vanished, the left side
would be identically zero, so this asymptotic proves the claim.

Statement A.4.8 gives the exact representation
\[
 Q_{j,I}=\int_{\mathbb R}u_j(i e^u)\,q_{t_j}(u)\,du,
 \qquad
 q_t(u)=
 A(u)e^{-it\tau(u)}
   \left(1+\sum_{n=1}^{N}t^{-n}b_n(u)+r_{N,t}(u)\right),
\]
where
\[
 A(u)=\frac{a}{2\cosh^{3/2}\tau(u)},\qquad
 r_{N,t}(u)=O_N(t^{-N-1}).
\]
The coefficients and the needed fixed-order derivatives are bounded
after the common factor $A(u)$ is removed.
In particular $A(u)\ll e^{-3|u|/2}$.
The Volterra recurrence in A.4.8 permits $N$ to be arbitrarily large.

Here is a precise way to preserve finite propagation in using this
expansion. Expand each $t^{-n}$ to sufficiently many Taylor terms
about $t=T$ under $\rho(t-T)$.
Every retained term is a polynomial in $t-T$ times $\rho(t-T)$,
multiplied by a power of $T^{-1}$.
Its negative-frequency reflection makes it an even entire
Paley--Wiener multiplier. For the leading term this is
\[
 h_{T,u}(t)=
 \rho(t-T)e^{it(d_0-\tau(u))}
 +\rho(-t-T)e^{-it(d_0-\tau(u))}.
\]
On $t>0$ the second term is rapidly small in $T$.
Its Fourier transform is supported in the two intervals of radius
$\varepsilon$ about $\pm(d_0-\tau(u))$.
Polynomial factors merely differentiate that compactly supported
Fourier transform.
All discarded spectral terms can be bounded before geometric
unfolding: use the compact-source eigenfunction sup bound,
the polynomial individual target-period bound, Weyl's law, the
rapid decay of $\rho$, and $\int A(u)\,du<\infty$.
Choosing the number of expansion terms larger than the resulting
fixed polynomial loss makes their total $O(T^{-2})$.
This step does not assert finite propagation for a gamma quotient.

\emph{2. The point-to-geodesic spectral density on the cover.}
We record the local kernel, including its normalization.
In Fermi coordinates $(r,v)$ about a complete geodesic in
$\mathbb H^2$, the transverse measure is $\cosh r\,dr$.
The positive Laplacian acting on functions independent of $v$
is $-\partial_r^2-\tanh r\,\partial_r$.
Conjugation by $\sqrt{\cosh r}$, followed by subtraction of $1/4$,
gives $-\partial_r^2+\operatorname{sech}^2r/4\geq0$ on the full line.
There are no negative bound states; a zero bound state is excluded
by its quadratic form, and positive bound states are excluded by
the nondecaying incoming/outgoing asymptotics at infinity.
The explicit Jost resolvent has smooth boundary values for $t>0$,
so the following jump is the full spectral density of the line
source, with no omitted discrete transverse spectrum.
The resolvent of its line delta source is
\[
 G_s(r)=-\frac{f_s(|r|)}{2f_s'(0)},\qquad \Re s>1.
\]
Indeed it solves the homogeneous equation away from zero and has
derivative jump $-1$, as required by the positive Laplacian.
Its outgoing boundary values at $s=1/2\mp it$ determine the
spectral density in $t$ by the resolvent jump formula:
\begin{equation}\label{A:line-spectral-density}
 \mathcal P_t(r)=
 \frac{t}{\pi i}
 \left\{-\frac{f_{1/2-it}(|r|)}{2f_{1/2-it}'(0)}
       +\frac{f_{1/2+it}(|r|)}{2f_{1/2+it}'(0)}\right\}.
\end{equation}
The factor is fixed by the scalar identity
$(x-\lambda-i0)^{-1}-(x-\lambda+i0)^{-1}
 =2\pi i\,\delta(x-\lambda)$ in the spectral resolution,
and by $d\lambda=2t\,dt$ for $\lambda=1/4+t^2$.
The denominators do not vanish for $t>0$: otherwise the
Wronskian of the incoming and outgoing solutions would vanish
at zero, whereas its value at infinity is nonzero.
The outgoing Volterra calculation of A.4.8 applies also on $r\geq0$.
It yields
\begin{equation}\label{A:line-spectral-principal}
 \mathcal P_t(r)=
 \frac{e^{itr}+e^{-itr}}{2\pi\sqrt{\cosh r}}
 +\frac1{\sqrt{\cosh r}}
       \sum_{\sigma=\pm1}e^{i\sigma tr}
           \sum_{n=1}^{N}t^{-n}c_{\sigma,n}(r)
 +\mathcal R_N(t,r).
\end{equation}
The coefficients and their fixed-order derivatives are uniformly
bounded; after removing their limits at infinity their derivatives
decrease exponentially.
For every fixed $N$ the remainder can be made
$O_N(t^{-N-1}\cosh^{-1/2}r)$ by taking sufficiently many Volterra
terms. This is uniform for $r\geq0$.
Only $r\geq d_0-\varepsilon>0$ will contribute below.

The normalization can also be checked without an asymptotic formula.
The square-root connection coefficient of the hypergeometric function
at argument one gives
\[
 f_s'(0)=-\frac{2^{1-s}\sqrt\pi\,\Gamma(s+1/2)}
                         {\Gamma(s/2)^2}.
\]
Consequently the source kernel is exactly
\[
 q_t(u)=\frac{a}{2\cosh\tau(u)}
              \frac{f_{1/2+it}'(\tau(u))}{f_{1/2+it}'(0)}
        =-\frac{a}{\cosh\tau(u)}G_{1/2+it}'(\tau(u)).
\]
This identifies the source factor with the normal derivative of
the same outgoing line resolvent. In particular no unspecified
scattering phase remains in its positive principal amplitude.

For an even Paley--Wiener multiplier, the automorphic point-to-cycle
kernel is obtained by summing this complete-line kernel over
the translates of the two target axes.
This follows either by periodizing the wave kernel and integrating
over a primitive target period, or directly by unfolding that integral.
Finite propagation makes the geometric sum finite at each fixed
source point and fixed propagation time.
The construction is valid on our compact orbifold; equivalently,
one can perform it on the fixed torsion-free cover and project to
the invariant subspace with the normalization factors already
verified in A.3.1.
Unoriented axes can occur with fixed positive multiplicities from
the cyclic stabilizer. We retain these multiplicities throughout;
their positivity is all that is needed here.

We include the finite-propagation justification relevant to the
shifted operator. For a smooth solution of
$w_{vv}+(\Delta-1/4)w=0$ with compactly supported initial data,
use the local energy density
$\mathcal E=(|w_v|^2+|\nabla w|^2+|w|^2)/2$.
Its conservation identity has an additional term
$5\operatorname{Re}(\overline w\,w_v)/4$, bounded by
$5\mathcal E/4$.
Integrate over the complement of the expanding unit-speed
neighborhood of the initial support.
The moving-boundary flux is nonpositive by
$|\operatorname{Re}(\overline{w_v}\partial_n w)|
 \leq(|w_v|^2+|\partial_nw|^2)/2$.
The exterior energy therefore satisfies
$E'(v)\leq5E(v)/4$ and $E(0)=0$.
Gronwall's inequality gives zero exterior energy.
Approximation extends the support assertion to distributional
initial data and hence to the wave kernel.
The argument is local, so it holds on the fixed smooth cover
and descends to the orbifold invariant subspace.
Fourier inversion of an even Paley--Wiener multiplier expresses
it through this cosine wave kernel, proving the support assertion
used in unfolding.

These rearrangements remain absolutely integrable in $u$.
Indeed, since the source is closed, the number of target axes at
distance at most $D$ from $p(u)$ is $O(e^D)$ uniformly in $u$.
To see this, translate the nearest point on each axis into a fixed
primitive segment of length $L_i$. A fixed basepoint of that segment
then lies within $D+L_i$ of $p(u)$. The resulting orbit points have
uniformly bounded multiplicity and can be bounded by disjoint-ball
packing. The source parameter may be reduced modulo its fixed period,
so the constants are uniform in $u$.
The propagated kernel vanishes unless
$r\leq\tau(u)-d_0+\varepsilon$.
The amplitudes in \eqref{A:line-spectral-principal} therefore give
\[
 \sum_{\gamma:r_\gamma\leq\tau(u)+O(1)}
       (\cosh r_\gamma)^{-1/2}
       \ll e^{\tau(u)/2}.
\]
Multiplication by $A(u)\ll e^{-3\tau(u)/2}$ leaves an integrable
$O(e^{-\tau(u)})$ bound.
A polynomial in $\tau(u)$, arising in an auxiliary remainder
estimate, is harmless as well.

We make explicit the support of the integrated spectral remainder.
The exact $\mathcal P_t(r)$ is even in $t$.
After reflection to the full real line, expand its high-frequency
symbol about $T$ into polynomials in $t-T$ under $\rho(t-T)$.
This Taylor construction and even reflection are applied to
\emph{every} inverse power $t^{-n}$ appearing in the cover
expansion \eqref{A:line-spectral-principal}, as well as to every
such power in the source expansion of $q_t$.
We never apply finite propagation directly to a retained
nonentire factor $t^{-n}$.
Every retained term has Fourier support in
$|r-|d_0-\tau(u)||\leq\varepsilon$ (with its reflected branch);
in particular it vanishes for
$r>|d_0-\tau(u)|+\varepsilon$.
The exact even Paley--Wiener kernel vanishes there by finite
propagation as well. Hence the integrated remainder also vanishes
there, and its uniform
$O_N(T^{-N})\cosh^{-1/2}r$ bound is summed only over this finite
geometric set. The compact low-$t$ part satisfies the same bound:
the displayed gamma formula makes $f_s'(0)$ nonzero also at $t=0$,
and the explicit hypergeometric formula bounds
$\sqrt{\cosh r}\,f_{1/2+it}(r)$ uniformly for bounded $t$.
The rapid decay of $\rho(t-T)$ supplies every inverse power of $T$.
Thus no unrestricted image sum of an unsmoothed spectral remainder
has been interchanged.

The constant eigenfunction contributes only
$O_N(T^{-N})e^{\tau(u)/2}$ before this multiplication:
evaluate the even entire multiplier at $t=i/2$ and use the
rapid horizontal decay of $\rho$ in a fixed strip.
No nonconstant exceptional eigenvalues exist by A.2.5.

\emph{3. The geometric oscillatory integral.}
Let $r_\gamma(u)\geq0$ denote the distance from $p(u)$ to a
target translate, counted with the fixed positive multiplicity.
For the leading terms, Fourier inversion in the window is exact.
The branch with phase $d_0-\tau-r_\gamma$ vanishes because
$\tau\geq R=2d_0$ and $\varepsilon<d_0$.
The other branch gives
\begin{equation}\label{A:wave-leading-integral}
 \frac1{2\pi}\sum_\gamma
 \int_{\mathbb R}
 \frac{a\,\widehat\rho(\tau-r_\gamma-d_0)}
      {2\cosh^{3/2}\tau\,\cosh^{1/2}r_\gamma}
       e^{iT(d_0-\tau+r_\gamma)}\,du.
\end{equation}
Choose a compact interval
$J\subset(0,R)$ containing $d_0$ and the support in
$\ell=\tau-r_\gamma$ of this expression.
All subleading terms have the same form, with smooth bounded
amplitude factors and extra powers of $T^{-1}$, and derivatives
of $\widehat\rho$ from the polynomial Taylor factors.
These derivatives have the same compact support.
Their bounds will follow from the same pushforward and stationary
phase arguments as the leading term.

\emph{4. A summable endpoint density, without a spacing assumption.}
Use unit Lorentz normals
\[
 n=(N_0,N_1,N_2)
   =\left(\frac{s}{\sqrt2},\sqrt{\frac32}\,t,3z\right),
 \qquad N_0^2-N_1^2-N_2^2=-1.
\]
On each choice of sign making $r_\gamma>0$, put $v=\tanh u$ and
\[
 D(v)=\sqrt{a^2-v^2},\quad
 E(v,\ell)=D(v)\cosh\ell-B\sinh\ell.
\]
The equation $\ell=\tau-r_\gamma$ becomes
\begin{equation}\label{A:wave-phase-equation}
 F(v,\ell)=N_0-N_1v-B\cosh\ell+D(v)\sinh\ell=0.
\end{equation}
On $J\subset(0,R)$, both $E$ and
$B\cosh\ell-D\sinh\ell$ are bounded away from zero.
The principal measure has the exact simplification
\[
 \frac{a}{2\cosh^{3/2}\tau\cosh^{1/2}r_\gamma}\,du
       =\frac{a}{2D^{3/2}E^{1/2}}\,dv .
\]
Moreover
\[
 F_v=-N_1-\frac{v\sinh\ell}{D},\qquad
 F_\ell=E,\qquad
 \frac{dv}{d\ell}
     =\frac{E}{N_1+v\sinh\ell/D}.
\]
If $|N_1|$ exceeds a fixed constant depending only on $J$,
the map has no critical point on this phase window.
Its pushed-forward density, including the smooth window, is
\begin{equation}\label{A:wave-pushforward-density}
 h_n(\ell)=
 \widehat\rho(\ell-d_0)
 \frac{a\sqrt{E(v(\ell),\ell)}}
      {2D(v(\ell))^{3/2}
       |N_1+v(\ell)\sinh\ell/D(v(\ell))|}.
\end{equation}
It exists on an interval, possibly empty. Extending it by zero,
including the possible endpoint jumps, gives
\[
 \|h_n\|_{L^1}+\operatorname{Var}(h_n)\ll_J |N_1|^{-1}.
\]
For completeness, $|N_1+v\sinh\ell/D|\geq |N_1|/2$,
$|v'|\ll_J|N_1|^{-1}$, $D$ and $E$ are bounded above and below,
and direct differentiation of \eqref{A:wave-pushforward-density}
gives $\sup|h_n|+\sup|h_n'|\ll_J|N_1|^{-1}$.
There are at most two endpoint jumps.
The allowed interval follows from
$F(\pm1,\ell)=N_0\mp N_1-Be^{-\ell}$.
This proves the stated BV estimate without separating any
two endpoint Busemann constants.

This bound is summable over all but finitely many target translates.
For $|N_2|>1$ orient the normal by $N_0>0$ and put
$c=N_0-|N_1|$.
The phase equation forces $0<c\leq C_J$.
The normal relation gives
\[
 |N_1|=\frac{N_2^2-1-c^2}{2c},
 \qquad \|h_n\|_{BV}\ll_J\frac{c}{N_2^2}
\]
for sufficiently large $|N_2|$.
For a fixed reduced target axis, its source-stabilizer translates
have
\[
 c_m=\sqrt{N_2^2-1}\,e^{-|u_0+mL_0|},
 \qquad
 \sum_{m:c_m\leq C_J}c_m
       \leq\frac{2C_J}{1-e^{-L_0}}.
\]
Statement A.3.1 gives $O(Y)$ reduced axes with $|N_2|\leq Y$.
Thus dyadic summation of $N_2^{-2}$ proves convergence of the
total BV norm at large height.
There are finitely many reduced axes at bounded height, including
the intersecting axes with $N_2=0$.
For each, $|N_1(m)|\asymp e^{|m|L_0}$ in its two infinite tails;
a zero leading coefficient would mean a shared endpoint with
the source, excluded by discreteness and the distinct axes.
Their inverse $|N_1(m)|$ sums converge too.
Only finitely many actual translates with bounded $|N_1|$ remain.

Consequently all infinite tails of \eqref{A:wave-leading-integral}
are Fourier transforms of a compactly supported BV density with
finite total variation. Their contribution is $O(T^{-1})$.
The same proof applies to the subleading amplitudes:
on the phase surface $\cosh\tau=D/\sqrt{1-v^2}$ and
$\cosh r_\gamma=E/\sqrt{1-v^2}$.
The outgoing amplitude coefficients, and their first derivatives
in $(v,\ell)$, are uniformly bounded up to $v=\pm1$.
This follows from the exponentially decreasing derivatives of
the Volterra coefficients; the factor $(1-v^2)^{-1}$ from
differentiating $\tau$ is canceled by their $e^{-2\tau}$ decay.
The BV product estimate therefore applies term by term.

\emph{5. All remaining critical values and Hessian signs.}
Only finitely many actual translates remain.
At an interior critical point of $\ell(v)$,
\eqref{A:wave-phase-equation} and $F_v=0$ give
\[
 N_1=-\frac{v\sinh\ell}{D},\qquad
 N_0=B\cosh\ell-\frac{a^2\sinh\ell}{D}.
\]
Substitution in the normal equation gives
\[
 N_2=\pm a\left(\cosh\ell-\frac{B\sinh\ell}{D}\right).
\]
The expression in parentheses is positive for $\ell\in J$.
Choose
$n'=(B,0,\operatorname{sgn}(N_2)a)$.
It is the unit normal of one of the translates $U^{\pm1}g_0$,
and direct substitution gives
\[
 \langle n,n'\rangle=-\cosh\ell.
\]
In the original integral coordinates this inner product is
$12s\mp51z$, a multiple of $3$.
Its absolute value is $3|z'|$ for one of the integral transformed
coordinates $z'=17z\mp4s$.
It follows that every positive critical value satisfies
\[
 \cosh\ell=3|z'|,\qquad |z'|\geq1.
\]
Choose $J$ small enough around $d_0$ that
$\ell<\operatorname{arcosh}(6)$ on $J$.
Then every critical value in this window equals $d_0$.

All such critical points are nondegenerate and have the same
Hessian sign. Indeed
\[
 F_{vv}=-\frac{a^2\sinh\ell}{D^3},\qquad
 \ell''(v)=-\frac{F_{vv}}{F_\ell}
     =\frac{a^2\sinh\ell}{D^3E}>0.
\]
Thus the phase $d_0-\ell(v)$ always has negative second derivative.
There is at least one such point: the target normal
$(2\sqrt2,0,3)$, corresponding to $(s,t,z)=(4,0,1)$,
has $v=0$, $\ell=d_0$, and
\[
 E=3,\qquad \ell''(0)=\frac{2\sqrt2}{51}.
\]
Its principal amplitude is strictly positive.
In fact the stationary set is explicit. At $\ell=d_0$,
$|N_1|\leq\sinh(d_0)/B=1/6$.
Since $N_1=\sqrt{3/2}\,t$ with integral $t$, this forces $t=0$,
then $v=0$, $N_0=2\sqrt2$, and $N_2=\pm3$.
Thus the only two geometric axes are those with normals
$(4,0,\pm1)$ in the original coordinates, each with its positive
cyclic-stabilizer multiplicity.
For one unit of such multiplicity the coefficient of
$e^{-i\pi/4}T^{-1/2}$ is
$\widehat\rho(0)/(4\,2^{1/4}\sqrt\pi)$.

There is no endpoint of any of these finitely many branches
with $\ell=d_0$.
At $v=\pm1$ the phase equation would require
$N_0\mp N_1=B e^{-d_0}$, hence
\[
 s\mp\sqrt3\,t=72-48\sqrt2.
\]
The left side belongs to $\mathbb Q(\sqrt3)$ and the right side
is a nonrational element of $\mathbb Q(\sqrt2)$, which is impossible.
We can therefore shrink the fixed $\varepsilon$ further so that
the finite-translate windows are compactly supported away from
their endpoints.
Ordinary stationary phase applies to their finitely many critical
points, as can be seen directly in one dimension.
Near a critical point $v_0$ set
$y=\operatorname{sgn}(v-v_0)\sqrt{2(\ell(v)-d_0)}$.
Then $dy/dv|_{v_0}=\sqrt{\ell''(v_0)}>0$ and the phase becomes
$-y^2/2$.
The leading integral is the amplitude at zero multiplied by
\[
 \int_{\mathbb R}e^{-iTy^2/2}\,dy
       =\sqrt{\frac{2\pi}{T}}e^{-i\pi/4}.
\]
This identity follows by inserting $e^{-\eta y^2}$, evaluating
the ordinary Gaussian integral, and letting $\eta\downarrow0$.
Taylor expansion of the smooth amplitude gives the usual
$O(T^{-3/2})$ local remainder; outside these neighborhoods,
integration by parts with $(iT\Phi')^{-1}\partial_v$ gives
$O(T^{-1})$.
Every leading coefficient is a positive real multiple
of the same factor $e^{-i\pi/4}T^{-1/2}$.
The value $\widehat\rho(0)>0$, the positive axis multiplicities,
and the displayed example show that their sum is nonzero.
Nonstationary parts are $O(T^{-1})$, as are the infinite BV tails.
All subleading symbol terms have an additional inverse power of
$T$, and the pre-unfolding remainder was made $O(T^{-2})$.
This proves \eqref{A:nonvanishing-wave-asymptotic}.

Finally group \eqref{A:nonvanishing-window} by complete eigenspaces.
Its nonzero asymptotic proves that at least one $\mathcal C(t)$
is nonzero. In fact finitely many nonzero contractions would make
\eqref{A:nonvanishing-window} $O_N(T^{-N})$ for every $N$, by
the Schwartz decay of $\rho$; this contradicts the same asymptotic.
Thus infinitely many are nonzero.
Statement A.5.6 then gives both claimed signs of
square-root oscillation for the actual sharp count.
\end{proof}
